\chapter{System Architecture and Implementation}
\label{chap:implementation}

\section{Pipeline architecture}

The implementation separates source acquisition, identity/topic processing,
graph materialization, validation, offline analysis, and interactive exploration.
Python stages use explicit input, configuration, and output paths. The central
graph artifact is a versioned JSON object validated against Pydantic contracts;
small manifests, configurations, and source ledgers are tracked, while raw API
caches, the full graph JSON, and database volumes remain local. The pipeline
does not silently replace unavailable stages with empty outputs.

\begin{figure}[tbp]
  \centering
  \begingroup\hyphenpenalty=10000\relax\exhyphenpenalty=10000\relax
\begin{tikzpicture}[
  font=\small,yscale=0.85,
  box/.style={draw=UETBlue,fill=PaleBlue,rounded corners=3pt,align=center,
    text width=6.2cm,minimum height=1.4cm,inner sep=5pt},
  store/.style={box,draw=UETGold,fill=UETGold!8,text width=4.1cm,minimum height=1.8cm},
  wide/.style={box,text width=13.5cm,minimum height=1.1cm},
  layer/.style={font=\small\bfseries,color=Slate,anchor=west},
  arrow/.style={-{Latex[length=2.1mm]},semithick,color=UETBlue}
]
  \node[layer] at (-7.2,1.2) {OFFLINE CONSTRUCTION};
  \node[box] (input) at (-3.7,0) {\textbf{Source inputs}\\Frozen rosters and source cache\\Versioned configuration};
  \node[box] (pipeline) at (3.7,0) {\textbf{Python processing pipeline}\\Matching, identity and topics\\Graph building and validation};
  \draw[arrow] (input) -- (pipeline);
  \node[wide] (graph) at (0,-1.8) {\textbf{Versioned graph and evidence artifacts}\\
    Graph JSON, accepted edges, source evidence and author--paper--organization records};
  \draw[arrow] (pipeline.south) -- (pipeline.south |- graph.north);
  \node[store] (neo4j) at (-5,-4.15) {\textbf{Neo4j Community}\\Isolated release graph\\Typed graph queries};
  \node[store] (sqlite) at (0,-4.15) {\textbf{SQLite ledger}\\Evidence / affiliations\\Graph JSON index};
  \node[store] (analysis) at (5,-4.15) {\textbf{Offline analysis}\\Frozen projection metrics\\Report figure artifacts};
  \draw[arrow] (graph.south -| neo4j.north) -- (neo4j.north);
  \draw[arrow] (graph.south) -- (sqlite.north);
  \draw[arrow] (graph.south -| analysis.north) -- (analysis.north);
  \node[wide] (api) at (0,-6.65) {\textbf{Serving layer: FastAPI backend}\\
    Read-only graph/evidence queries and precomputed analytics};
  \draw[arrow] (neo4j.south) -- (neo4j.south |- api.north);
  \draw[arrow] (sqlite.south) -- (api.north);
  \draw[arrow,dashed] (analysis.south) -- (analysis.south |- api.north);
  \node[wide] (react) at (0,-8.95) {\textbf{Presentation layer: React / TypeScript Explorer}\\
    Sigma.js + Graphology: graph views\qquad ECharts: analytical views};
  \draw[arrow] (api.south) -- (react.north);
\end{tikzpicture}
\endgroup

  \caption{Implementation architecture for the frozen v4 release. Batch
  construction and analysis produce versioned local artifacts; the Explorer
  reads the isolated graph and evidence stores. The dashed route supplies
  precomputed analysis artifacts to the Explorer.}
  \label{fig:architecture}
\end{figure}

\section{Graph storage and integrity}

The released graph is loaded into the isolated Neo4j Community 5.26 service
\modelid{neo4j-five-year} \citep{neo4jDocs}. The loader checks the configured graph version and
uses bounded transactions. The release's recorded database verification
reconciles node and relationship counts, executes six integrity queries, and
compares a repeated import with the original per-label and per-relation maps.
Those records report 47,875 nodes and 225,731 edges after replay. The pilot and
extended databases are separate; they are not inputs to this report.

Graph JSON, analysis JSON, configuration hashes, content-freeze receipts, and
change ledgers bind the analysis to graph \GraphVersion{} and methods
\modelid{release-openalex-graph-v1.1} and
\modelid{named-projection-analysis-v2}. Evidence records retain source URLs,
retrieval times, raw hashes, source locators, and permitted excerpts. Derived
evidence points to accepted base edges and distinct core-paper support. The
database is a queryable deployment of the graph artifact; it is not the sole
record of the graph's meaning or provenance.

\section{Analysis implementation}

Offline analysis uses Python, NetworkX \citep{networkxDocs}, and Matplotlib
\citep{matplotlibDocs}. The configured analysis
produces separate projection metrics, annual counts, community assignments,
centrality rankings, topic shares, sector/country mixing, and CSV-backed figures.
The immutable release-v4 output records the graph and configuration hashes,
seeds, methods, and known limitations. Static figures in this report are copies
of these measured release artifacts; they were not regenerated from a new graph
for this manuscript.

\section{Research Explorer}

The Explorer is a read-only application for inspecting the same release. Its
backend uses FastAPI and Pydantic and issues parameterized read queries through
the Neo4j Python driver. A hybrid artifact store supplies precomputed
full-projection metrics, while an immutable SQLite ledger indexes evidence and
author--paper--organization observations. This arrangement keeps provenance
inspection available without requiring every request to materialize the full
graph ledger.

The frontend uses React and TypeScript. Sigma.js with Graphology renders
interactive node-link views; ECharts presents the precomputed temporal,
collaboration, citation, and quality summaries. The application supports graph
projection and filtering, entity/edge inspection, evidence navigation,
bounded neighborhood exploration, comparisons, tables, and saved views. These
features make the analysis easier to inspect; they do not independently validate
the underlying claims. The recorded project URLs and access caveat appear in
\cref{chap:reproduction}.

\section{Reproducibility and version discipline}

The source repository locks Python dependencies in \texttt{uv.lock}. Analyses
use named configuration versions and new output paths for replays. The v4 graph
was created as a separate correction from its predecessor, with protected
paper, citation, topic, taxonomy, and non-derived evidence equality checks. The
experimental v5 identity work is maintained separately and has not replaced v4.
Reproduction requires restoring permitted ignored graph and OpenAlex inputs and
checking their recorded hashes; acquiring current provider responses is not
byte-identical restoration. Appendix B provides the local commands and source
identifiers needed to repeat available checks.
